Die gesamte Stoffmenge bleibt erhalten. Bei gleicher Temperatur ist $n = pV/(RT)$, also
\begin{align}
 \frac{p_1V_1}{RT}+\frac{p_2V_2}{RT} &= \frac{p\left(V_1+V_2\right)}{RT}
\end{align}
(Man kann auch jede Gasportion für sich nach Boyle-Mariotte auf das Gesamtvolumen ausdehnen und die beiden Teildrücke addieren.) Der gemeinsame Druck ist
\begin{align}
 p &= \pMF \\
   &= \frac{\pEi\cdot\VEi+\pZw\cdot\VZw}{\VEi+\VZw} \\
   &= \pM \approx \result{\pMbP{0}{2}}
\end{align}
Das Volumen des Röhrchens ist vernachlässigt.
